\section{Joining the evidence}
\label{sec:interpretation}

The two studies define a staged systems test. Cohort supply determines whether
any legal scheduler can find enough same-group work before launch deadlines.
Observation placement determines whether an already resident decision chain
pays a host copy, synchronization, branch, and redispatch epoch. A proposed GPU
route service must pass both gates. Failure at either gate is a concrete reason
to keep that transition on the CPU.

The current numbers do not form a service-level acceleration estimate and must
not be multiplied.
The trace study sweeps candidate thresholds, while the resident study fixes
three cohort sizes and a synthetic route body. The trace also contains no
measurement that 32 consecutive control epochs can execute without a model,
tool, policy, or commit observation. The value of the separation is that each
missing quantity now has a named measurement rather than being hidden inside a
broad GPU-agent claim.

\subsection{The next systems test}

The runtime implied by the boundary is a finite-capacity online route compactor.
It receives typed completion events, updates compact state, queues events by a
verified executable route and deadline, launches a route only above its
measured safe-suffix threshold $K_r$, and falls back to a tuned CPU path when a
deadline or capacity limit prevents batching. The GPU may emit ordered effect
descriptors; a privileged CPU or DPU remains the authority for external
effects.

Such a system adds one essential quantity. Let $z_i=1$ only when the frozen
online policy accelerates event $i$ with an exact output and observed device
start no later than its launch deadline; otherwise $z_i=0$. Then

\begin{equation}
  A=\frac{1}{|E|}\sum_{i\in E}z_i.
\end{equation}

Late, missing, failed, and fallback events have $z_i=0$ and remain in $E$. The
opportunity recovery above fixed windows is

\begin{equation}
  R_A=\frac{A-F}{\pstar-F},
  \label{eq:online-recovery}
\end{equation}

when $\pstar>F$. This turns the offline result into a direct runtime target. A
credible evaluation pairs $A$ and $R_A$ with raw invocation P99, deadline
misses, CPU core-seconds per event, exact trajectories, task utility, network
bytes, cost, and, for a shared inference GPU, TTFT and TPOT guardrails.

\subsection{Deployment paths}

The strongest architecture is a regional route service that aggregates typed
events from many inference and tool workers. Aggregation attacks the cohort
supply problem directly; measured ingress, state movement, and effect egress
then determine whether the control-path savings survive the network boundary.
A second path uses a low-priority queue on an inference GPU that is already
allocated. Its economic advantage is marginal device cost, while its decisive
constraint is noninterference with model TTFT and TPOT.

A dedicated control GPU is the weakest starting point because the present
experiments do not establish enough utilization or CPU displacement to pay for
one. GPU-initiated VM lifecycle operations are also outside the transition
unit: a device can emit a typed prewarm request, but a privileged CPU or DPU
control plane must authorize and commit it.
